% !TEX root = ../main.tex

%----------------------------------------------------------------------------------------
% CHAPTER 2 - ESTUDI DE VIABILITAT
%----------------------------------------------------------------------------------------

\chapter{Estudi de viabilitat}
\label{Ch:Viability}

%----------------------------------------------------------------------------------------

\section{Viabilitat econòmica}

El projecte ha estat desenvolupat amb recursos mínims:

\begin{itemize}
    \item \textbf{Hardware}: Ordinador personal del desenvolupador (sense costos addicionals)
    \item \textbf{Software}: Tecnologies completament open-source i gratuïtes:
    \begin{itemize}
        \item Node.js i npm per a el pipeline ETL i frontend
        \item React per a la interfície web
        \item GitHub Actions per a l'automatització (gratuït per a repositoris públics)
    \end{itemize}
    \item \textbf{Dades}: APIs públiques gratuïtes de la Generalitat de Catalunya (Socrata)
    \item \textbf{Hosting}: Pot ser desplegat a GitHub Pages (gratuït) o a altres plataformes low-cost
\end{itemize}

En conclusió, el cost econòmic del projecte és negligible, ja que utilitza exclusivament recursos gratuïts i infraestructura oberta.

\subsection{Estimació de cost per hores de desenvolupament}

Encara que el cost directe de tecnologies i infraestructura és nul (per l'ús exclusiu de recursos gratuïts i open-source), el projecte requereix una inversió significativa en temps de desenvolupament humà. A continuació es presenta una estimació hipotètica del cost del projecte basada en hores de treball real:

\begin{table}[h]
\centering
\caption{Estimació de cost del projecte per fases (pla hipotètic)}
\begin{tabular}{|l|c|r|r|}
\hline
\textbf{Fase} & \textbf{Hores} & \textbf{Cost/hora (€)} & \textbf{Total (€)} \\
\hline
Investigació i configuració & 30 & 25 & 750 \\
\hline
Desenvolupament pipeline ETL & 60 & 25 & 1.500 \\
\hline
Frontend - Dashboards 1 i 2 & 67 & 25 & 1.675 \\
\hline
Frontend - Dashboards 3 i 4 & 49 & 25 & 1.225 \\
\hline
Automatització i desplegament & 52 & 25 & 1.300 \\
\hline
Testing i validació & 42 & 25 & 1.050 \\
\hline
Tutories i revisions & 12 & 30 & 360 \\
\hline
Reserva d'imprevistos & 18 & 25 & 450 \\
\hline
\textbf{TOTAL} & \textbf{330} & & \textbf{8.310} \\
\hline
\end{tabular}
\label{tab:cost-estimation}
\end{table}

L'estimació total de 330 hores inclou el desenvolupament funcional (258 hores), el testing i validació (42 hores), les reunions periòdiques amb el tutor acadèmic (12 hores), i una reserva de 18 hores per a imprevistos i tasques de coordinació. El cost per hora s'ha estimat en 25 €/h per a tasques de desenvolupament i 30 €/h per a tasques de coordinació i revisió, xifres coherents amb els costos mitjans del mercat laboral espanyol per a perfils de desenvolupador web júnior.

Aquesta estimació demostra que, tot i utilitzar tecnologies gratuïtes, el cost real del projecte resideix principalment en la inversió de temps humà, representant un valor afegit significatiu que justifica la seva continuïtat i escalabilitat.

\subsection{Model de negoci hipotètic}

Per tal de finançar el manteniment i l'escalabilitat del projecte a llarg termini, es podrien plantejar diversos models de negoci hipotètics:

\begin{itemize}
    \item \textbf{Model de subscripció freemium}: L'accés bàsic al dashboard seria gratuït (finançat per publicitat contextual o patrocini d'entitats públiques), mentre que funcionalitats avançades com l'exportació de dades, l'API programàtica o les alertes personalitzades requeririen una subscripció mensual (5-15 €/mes per a professionals).
    
    \item \textbf{Contractació institucional}: Vendre llicències d'ús a ajuntaments, consells comarcals o agències de l'aigua que necessitin monitoritzar els recursos hídrics del seu territori. El model de preus es basaria en el nombre d'embassaments monitoritzats i la freqüència d'actualització (500-2.000 €/any per a administracions locals).
    
    \item \textbf{Patrocini i col·laboració amb ONG}: Entitats com Aigües de Catalunya o Greenpeace podrien patrocinar el projecte a canvi de visibilitat i dades exclusives per a campanyes de conscienciació ambiental.
    
    \item \textbf{Servei de consultoria de dades}: Ofereix anàlisis personalitzades i informes periòdics sobre l'estat dels recursos hídrics a empreses del sector agrícola, energètic o turístic que depenen de l'aigua (200-500 € per informe).
\end{itemize}

Tanmateix, cal destacar que l'objectiu primordial del projecte és la divulgació pública i l'accés obert a la informació hídrica, per la qual cosa el model més alineat amb aquests valors seria un servei gratuït suportat per patrocini institucional i col·laboracions amb entitats del sector públic i de la societat civil.

%----------------------------------------------------------------------------------------

\section{Viabilitat tecnològica}

Les tecnologies necessàries estan totalment disponibles i ben establertes:

\subsection{Extracció de dades}

\begin{itemize}
    \item \textbf{APIs Socrata}: Les dades dels embassaments i precipitacions estan disponibles a través de les APIs públiques de la Generalitat de Catalunya. S'ha verificat que aquestes APIs són estables i proporcionen més de 200.000 registres en total (~87.000 d'embassaments i ~135.000 de precipitació).
    \item \textbf{Formats de dades}: Les dades es retornen en JSON, format estàndard i fàcil de processar.
\end{itemize}

\subsection{Processament i transformació}

\begin{itemize}
    \item \textbf{Node.js}: Runtime de JavaScript robust i escalable per a backends.
    \item \textbf{npm packages}: Mòduls madurs com \texttt{axios} (client HTTP), \texttt{fs} (sistema de fitxers) estan ben establerts.
\end{itemize}

\subsection{Visualització}

\begin{itemize}
    \item \textbf{React}: Framework frontend modern amb ecosistema madur.
    \item \textbf{Recharts}: Biblioteca de gràfics construïda sobre React, perfecta per a visualitzacions de sèries temporals i gràfics compostos.
    \item \textbf{Leaflet}: Biblioteca de mapes interactius de codi obert, ideal per a visualitzar embassaments en coordenades geogràfiques.
\end{itemize}

\subsection{Automatització}

\begin{itemize}
    \item \textbf{GitHub Actions}: Plataforma de CI/CD nativa de GitHub que permet executar workflows automàtics (ETL) en horaris programats.
    \item \textbf{Cron jobs}: Suportats de forma nativa per a programar execucions diàries.
\end{itemize}

\textbf{Conclusió}: Totes les tecnologies necessàries estan disponibles, ben documentades, i han estat provades exitosament durant el desenvolupament del projecte.

%----------------------------------------------------------------------------------------

\section{Viabilitat ètica i legal}

\subsection{Privacitat de dades}

Les dades utilitzades en aquest projecte provenen de fonts públiques de la Generalitat de Catalunya:
\begin{itemize}
    \item Els nivells dels embassaments són dades públiques de caràcter no personal.
    \item Les dades de precipitació són recopilades per estacions meteorològiques i són agregades (sense identificar ubicacions específiques).
    \item No es tracten dades personals de ciutadans.
\end{itemize}

\subsection{Compliment legal}

\begin{itemize}
    \item \textbf{Accés a dades públiques}: El projecte respecta les condicions d'ús de les APIs públiques de Socrata.
    \item \textbf{Llicència open-source}: El codi és alliberat sota llicència MIT, permetent ús i modificació.
    \item \textbf{RGPD}: No s'aplica, ja que no es tracten dades personals.
\end{itemize}

\subsection{Seguretat}

\begin{itemize}
    \item La pipeline ETL utilitza connexions HTTPS per a comunicar-se amb les APIs.
    \item No s'emmagatzemen credentials sensibles en el repositori (utilització de \texttt{.env}).
\end{itemize}

%----------------------------------------------------------------------------------------

\section{Recursos}

\subsection{Recursos humans}

\begin{itemize}
    \item Un desenvolupador (autor del projecte) responsable del disseny, implementació i documentació.
    \item Aproximadament 330 hores de desenvolupament (incloent investigació, implementació, testing, documentació i tutories), distribuïdes en 5 fases al llarg de 12 setmanes.
\end{itemize}

\subsection{Recursos materials}

\begin{itemize}
    \item Ordinador de desenvolupament (ja disponible)
    \item Connexió a Internet (ja disponible)
    \item Accés a GitHub per a control de versions i CI/CD
\end{itemize}

%----------------------------------------------------------------------------------------

\section{Conclusions de la viabilitat}

El projecte és completament viable pels següents motius:

\begin{enumerate}
    \item \textbf{Cost econòmic nul}: S'utilitzen exclusivament tecnologies gratuïtes i open-source.
    \item \textbf{Tecnologies madures}: Totes les tecnologies necessàries estan ben establertes i comprovades en entorns de producció.
    \item \textbf{Dades disponibles}: Les APIs públiques de la Generalitat proporcionen totes les dades necessàries.
    \item \textbf{Legal i ètica}: No hi ha conflictes legals ni ètics, ja que es tracta de dades públiques sense identificació personal.
    \item \textbf{Escala factible}: El volum de dades (200.000+ registres) es pot processar en poc de temps, i l'actualització diària és completament viable.
\end{enumerate}

En conclusió, el projecte de visualització de dades hídriques de Catalunya és viable de realitzar amb els recursos disponibles i ofereix valor significatiu per a la societat civil i els responsables de política hídrica.

%----------------------------------------------------------------------------------------
